\section{Summary of Contributions}
\label{sec:conclusion-summary}

This thesis extended Goswami's two-layer range filter construction in
three areas: a single-vector ERE backend that cuts metadata size by
$24\%$ (Chapter~\ref{chap:ere}); an adaptive bucket-search procedure
that drops $59$~bits of per-key metadata and speeds queries by
$13$--$30\times$ at LSM bucket sizes (Chapter~\ref{chap:ere}); a
truncation hash that saves $\log_2 \mathcal{L}$ bits per key on sparse
inputs (Chapter~\ref{chap:hash}); and Seg-ARE, a single-pass
segmentation that routes dense clusters to exact-mode ERE sub-filters
and sparse tails to the truncation hash (Chapter~\ref{chap:hybrid}).
The first three improvements bring the Goswami-line filter closer to
the lower bound on adversarial inputs. Seg-ARE goes below the bound on
distributions where the structure permits.

\section{Seg-ARE and Grafite as Variants of Elias--Fano Coding}
\label{sec:conclusion-ef-view}

A useful way to read the empirical results of
Chapter~\ref{chap:evaluation} is through the Elias--Fano (EF) primitive
that both Grafite and Seg-ARE rely on at the structural level.

Grafite, in the regime where the requested $\varepsilon$ falls below
the lossless threshold $u/n$, degenerates into a single EF encoding of
the raw sorted keys --- the wrapper substitutes
\texttt{ef\_sux\_vector(keys)} and the filter answers range queries by
predecessor probe. Below this threshold Grafite is an approximate
filter; above it, it is essentially EF on the whole universe.

Seg-ARE applies the same EF primitive, but only after segmenting the
input. Each dense cluster identified by the forward sweep is handed to
an exact-mode ERE sub-filter, which encodes the cluster's keys
locally; sparse keys outside any cluster fall through to the
truncation hash. The total size of a Seg-ARE filter is the sum of
per-cluster EF encodings plus the fallback's hash table --- versus
Grafite's single EF encoding over the whole universe.

This framing makes the win condition concrete: Seg-ARE beats Grafite
exactly when per-cluster EF is more compact than universe-wide EF,
i.e.\ when the data has dense regions separated by large gaps the
segmentation can isolate. Quantitatively, the bpk gap between the two
filters on a clustered input decomposes as
\[
    \log_2(U_{\text{global}}/n) \;-\;
    \overline{\log_2(2^{M_c}/n_c)},
\]
where $U_{\text{global}}$ is the spread of the whole input and
$2^{M_c}, n_c$ are the local window width and key count of the $c$-th
cluster. On SOSD-Wikipedia the segmenter resolves the $\sim\!10^6$
unique timestamps into roughly a dozen epoch-aligned clusters, each
occupying a window of $2^{M_c} \leq 2^{25}$ around a $2^{26}$-wide
global spread; the per-key cost drops from $\log_2(U/n) \approx 6.1$
bits to $\overline{\log_2(2^{M_c}/n_c)} \approx 3.6$ bits, and the
difference shows up directly in the headline numbers above.

\section{When Seg-ARE Wins, Ties, and Loses}
\label{sec:conclusion-where}

The empirical picture across the four SOSD distributions falls out of
the EF interpretation above:

\begin{itemize}
    \item \textbf{Wiki --- clearest win.} Wikipedia edit timestamps
    are heavily concentrated on the right side of the universe
    (Figure~\ref{fig:sosd-histograms}), with dense regions and large
    sparse gaps in between. The segmentation isolates the dense parts
    cleanly; per-cluster EF encodes each tight region in far fewer
    bits than a universe-wide EF could. Seg-ARE reaches
    $\mathrm{FPR} < 10^{-8}$ at $\sim\!6$ BPK against Grafite's
    $\sim\!7$ BPK at the lossless threshold.
    Honest disclosure: on Wiki the segmenter alone is doing the work
    --- the cluster geometry routes nearly all empty queries to
    exact-mode sub-filters, and the truncation-hash fallback rarely
    fires. On Books and OSM, where the cluster/gap separation is
    weaker, the fallback path contributes a larger share of the
    overall accuracy.
    \item \textbf{Facebook --- tied.} Facebook user IDs occupy a
    relatively small $37$-bit universe with no large gaps. There is
    little for segmentation to exploit, and universe-wide EF is
    already compact. Seg-ARE and Grafite both reach FPR $= 0$ at
    $\sim\!11$ BPK; Seg-ARE matches but does not improve over the
    Grafite-with-fallback baseline.
    \item \textbf{OSM --- no advantage.} OpenStreetMap cell IDs are
    spread thinly across the full $64$-bit universe, with median gaps
    of order $10^7$. The DBSCAN-style detector finds essentially no
    clusters above the $\delta = 256\mathcal{L}/\varepsilon$ window,
    so most keys go to the truncation-hash fallback and Seg-ARE
    degrades to truncation. Honest hashing is the only option on this
    geometry; no segmentation primitive helps.
    \item \textbf{Uniform --- degrades to fallback.} Synthetic uniform
    distributions are the limit case of OSM: no clusters at any scale,
    every key in the fallback, and Seg-ARE collapses into plain
    truncation. This is the worst-case axis where the worst-case
    Goswami bound bites; here the contributions from
    Chapter~\ref{chap:ere} (single-vector ERE, adaptive bucket search)
    are what keep the filter close to the bound.
\end{itemize}

Seg-ARE is therefore not a universal improvement over Grafite. It is a
distribution-aware option that activates when the data has the
cluster-and-gap geometry common in many real LSM key sets
(Wiki-like, Facebook-trace-style burst patterns) and stays harmless on
the rest.

\section{Limitations}
\label{sec:conclusion-limits}

The Seg-ARE design carries three honest limitations.

\textbf{Detector sensitivity.} The forward sweep depends on
$\delta = \texttt{minPts} \cdot \mathcal{L} / (2\varepsilon)$ and on
the fixed choice $\texttt{minPts} = 512$
(Chapter~\ref{chap:hybrid}). On clusters smaller than $\texttt{minPts}$
keys the detector misses real density and the win is lost. We have
not investigated adaptive $\texttt{minPts}$ tuning.

\textbf{Per-segment overhead.} Each segment carries $\sim\!700$~bits
of fixed metadata (cluster bounds, sub-filter header, pointer). At
$\texttt{minPts} = 512$ this is $\sim\!1.4$~BPK of fixed overhead. On
inputs that produce many small clusters the overhead can outweigh the
per-cluster EF savings.

\textbf{Static setting.} The construction takes a fixed key set;
segments are not maintained under inserts or deletes. Dynamic key
sets, where the cluster structure itself evolves, require the
keepsake-box machinery of Memento~\citep{eslami2024memento} or an
equivalent online segmentation scheme.

\section{Future Work}
\label{sec:conclusion-future}

Three directions stand out.

\textbf{Adaptive segmentation parameters.} The window and density
thresholds are currently fixed at construction time from a single
\texttt{minPts} value. Choosing them per LSM level (top levels have
fewer keys but tighter clusters; bottom levels are denser but flatter)
could close the gap on OSM-like distributions where the global
$\delta$ is too coarse.

\textbf{Combining segmentation with learned models.} The CDF-fitting
approach of SNARF~\citep{vaidya2022snarf} solves the universe-mapping
problem that hurts Grafite on small-universe inputs, but loses the
worst-case guarantee. Running a segmenter first, then a CDF model on
each cluster, would localise the model's assumptions and might recover
the guarantee on the clusters individually.

\textbf{Dynamic Seg-ARE.} Combining the segmentation primitive with
Memento's dynamic keepsake-box machinery would extend the
distribution-aware win condition from static SSTables to mutable
indexes (B-trees, single-global-filter LSM-trees, mixed
transactional/analytical workloads).
